\section{接口、数据与生态}
\label{sec:interfaces}

\subsection{Facade：唯一的任务级入口}

所有能力收敛到 \texttt{Facade} 一个类（ADR 0014），CLI 与 MCP 服务由
Facade 方法元数据同源派生——任务级 API 模型除执行输入校验外，还公开
当前任务上下文下的参数元数据（条件取值域与校验器共用同一份规则定义），
供 GUI、CLI、MCP 生成控件与范围提示，外部消费者不维护本地参数规则
副本。任务级方法分三档：

\begin{itemize}
  \item \textbf{一档任务}：\texttt{design\_orbit}（任务轨道设计：
        CR3BP 初猜 $\to$ 星历修正 $\to$ 高精度预报）、
        \texttt{control\_orbit}（轨道保持 Monte Carlo 仿真）、
        \texttt{transfer\_design}（HMN/LGA/WSB/低推力）、
        \texttt{orbit\_propagation}（轨道预报），以及
        \texttt{spacetime\_transform}（会合系与 J2000、GCRS 与
        EBCRS 的时空互转）；
  \item \textbf{二档子任务}：\texttt{orbit\_family\_generation}
        （八族延拓生成，族记录自动入库）等；
  \item \textbf{轨道库}：7 个 catalog 方法（见下）。
\end{itemize}

\begin{lstlisting}[caption={设计一条地月 L2 Halo 轨道（examples/main\_design.py）},label=lst:quickstart]
from e2m2e.api import Facade

facade = Facade()

result = facade.design_orbit(
    orbit_type="Halo",
    collinear_point=2,
    amplitude=30000.0,          # km
    epoch=[2024, 1, 1, 0, 0, 0.0],
    duration=365.25 * 86400.0,  # s
)

print(result.orbit_type)
print(result.initial_state)
\end{lstlisting}

\subsection{MCP：Agent 原生接口}

e2m2e 可作为 MCP 服务器（\texttt{e2m2e mcp-serve}，stdio 传输，不监听
端口）向 LLM Agent 暴露 13 个任务级工具：6 个任务方法（一档 5 个 +
族生成）与 7 个轨道库工具。工具清单由 Facade 方法元数据纯派生，无需
另行维护。所有响应走统一信封 \texttt{\{status, data, error, meta\}}：
错误以错误码表达（\texttt{INVALID\_PARAMS}、\texttt{RECORD\_NOT\_FOUND}、
\texttt{PROPAGATION\_FAILED} 等），不泄漏 traceback；数值工具以
\texttt{data.status} 报告 converged / diverged / stagnated /
max\_iterations 软失败语义（ADR 0020）——Agent 可以把“未收敛”当作
正常信息来处理，而不是一次异常。

配置方式是在 MCP 客户端注册一条命令；建议把 \texttt{cwd} 钉在含
\texttt{kernels/}（SPICE 内核）与 \texttt{catalog/}（轨道库）的目录，
或以 \texttt{SPICE\_KERNEL\_DIR} / \texttt{E2M2E\_CATALOG\_DIR} 指定
绝对路径。配置完成后即可用自然语言驱动，例如：

\begin{quote}\itshape
设计一条近月点高度 3000 km 的 L2 南族 NRHO，然后对它做 100 次
Monte Carlo 轨道保持仿真，结果打上 ``candidate'' 标签。
\end{quote}

\subsection{轨道库 catalog}

计算产物自动归档（ADR 0031）：记录 = JSON 元数据 + NPZ 数组段（带
schema 版本号），SQLite 仅作派生索引。产物型工具成功后返回
\texttt{record\_id}，可跨工具链式引用（例如 \texttt{control\_orbit}
以 \texttt{input\_record\_id} 引用 \texttt{design\_orbit} 的产物）。
7 个库工具覆盖完整生命周期：\texttt{catalog\_query}（按族 / 平动点 /
Jacobi 区间 / 振幅 / 标签 / 收敛状态多维过滤）、\texttt{catalog\_get}、
\texttt{catalog\_delete}、\texttt{catalog\_tag}（教学标注）、
\texttt{catalog\_promote}（族成员提升为独立记录）、
\texttt{catalog\_export}（打包导出）、\texttt{catalog\_sweep}
（参数空间批量扫描入库）。

\subsection{安装与部署}

\begin{lstlisting}[style=e2bash,caption={安装与 SPICE 内核配置},label=lst:install]
uv pip install e2m2e            # 发布 wheel 内嵌 CSPICE，开箱即用
uv pip install "e2m2e[mcp]"     # MCP 服务

# 源码开发（需 Rust 工具链）：
git clone https://github.com/cislunarspace/e2m2e.git
cd e2m2e
make dev      # 同步依赖 + 拉取 CSPICE 编译包与 SPICE 内核 + 构建 Rust 扩展
\end{lstlisting}

发布 wheel 覆盖 Windows x86\_64、Linux x86\_64 与 Linux aarch64
（鲲鹏 / 飞腾 / 树莓派），内嵌静态链接的 CSPICE，Rust 快速路径
（STM 传播、打靶、第三体引力等）开箱即用；扩展不可用时显式报错而非
静默降级。SPICE 内核打包为 GitHub Release 资产 \texttt{kernels-v1}
（\texttt{make setup} 自动下载），包含 DE 行星星历、地球自转、月球
姿态、闰秒与行星常数内核；r2s2 时空转换另需含 TT$-$TDB 时间星历的
历表。官方来源为 NASA NAIF。

\texttt{examples/} 目录提供六个可运行示例：端到端轨道设计、轨道保持
Monte Carlo、高保真传播、Lambert 转移、环月轨道四层对比（LLO/ELFO/
DRO/Halo，全量运行不到 1 分钟）、轨道族生成。
